% RESEARCH_DRAFT_SHA256: 2456b2df59c1a85def19a6354e97bf4e938cb9441f1be97c851830ab8660b277
% Layout and source-reference wrappers expanded for the web renderer.
\documentclass[11pt]{article}
\usepackage[margin=25mm]{geometry}
\usepackage{fontspec,unicode-math}
\setmainfont{Latin Modern Roman}
\setmathfont{Latin Modern Math}
\usepackage{amsmath,amssymb,mathtools}
\usepackage{xurl,hyperref}
\hypersetup{colorlinks=true,urlcolor=blue,linkcolor=blue}
\setcounter{secnumdepth}{0}
\setlength{\parindent}{0pt}
\setlength{\parskip}{5pt}
\setlength{\emergencystretch}{3em}
\newcommand{\sref}[2]{\hyperlink{source:#1:#2}{[S#2]}}
\newcommand{\cataloguescope}[1]{\paragraph{Source question.}#1}
\newcommand{\floor}[1]{\left\lfloor#1\right\rfloor}
\begin{document}
\section*{Research-attempt provenance}
This standalone research attempt retains the complete original definition body below,
including its historical source and preparation notes. Those notes describe the
earlier catalogue preparation; the new attempt and its verification are identified
separately after that body. No original ID, catalogue status, or source statement is changed.

% UnsolvedMath ID 30000155; source catalogue code OWR-768-005
\section{Asymptotic Goodness of Function-Field-Tower Codes}\label{30000155}
{\small UnsolvedMath ID 30000155 · Computer Science}
\subsection{Definitions and mathematical statement}

\paragraph{Shared notation.}Standard number systems are denoted by $\mathbb N,\mathbb Z,\mathbb Q,\mathbb R,\mathbb C$; index ranges, computational models and specialized notation follow the cited problem source and the definitions below.


For a sequence of linear $[N_j,K_j,D_j]_q$ codes with $N_j\to\infty$, asymptotic goodness means $\liminf K_j/N_j>0$ and $\liminf D_j/N_j>0$. Algebraic-geometric codes evaluate suitable functions on rational points of curves over finite fields.

A code $C\subseteq\Sigma^N$ is a set of permitted length-$N$ words. Its Hamming minimum distance is the minimum number of unequal coordinates of two distinct words. For a linear $[N,K,D]_q$ code, rate is $K/N$ and relative distance is $D/N$; decoding recovers a codeword from corrupted observations.
\paragraph{Statement recorded in the supplied source.}
Are the algebraic-geometric code families obtained by the direct-sum construction from the first levels of recursively defined asymptotically good towers of function fields themselves asymptotically good, with positive asymptotic rate and positive relative minimum distance? \href{https://doi.org/10.4171/owr/2004/54}{[S2]}
\paragraph{Scope qualification.}
The tower, selected levels, and direct-sum construction are report-specific and are not reconstructed from the short supplied question.
\paragraph{Source status.}
The supplied catalogue labels this record open. Its source-status metadata was checked against the supplied snapshot on 2026-10-11. This import does not independently establish current open status or mathematical correctness. The catalogue literature assessment is dated 2026-08-21.
\subsection{Short English statement}
Determine whether the stated direct-sum code families retain positive rate and relative distance as their lengths grow.
\subsection{Sources}
\begin{itemize}
\item[S1] Ulam AI and the UnsolvedMath Contributors, UnsolvedMath: A Curated Collection of Open Mathematics Problems, record 30000155 (OWR-768-005). Catalogue curation and metadata: CC BY 4.0. The supplied source question is retained; definitions and the English summary are editorial additions.
\url{https://huggingface.co/datasets/ulamai/UnsolvedMath}.
\item[S2] Finite Fields: Theory and Applications (2005), Oberwolfach Reports, DOI 10.4171/owr/2004/54, pp. 2913--2970
\url{https://doi.org/10.4171/owr/2004/54}.
\item[S3] EMS article overview
\url{https://ems.press/journals/owr/articles/768}.
\textit{Underlying Oberwolfach Reports material retains its source-specific licence, including CC BY-SA 4.0 where applicable; see THIRD\_PARTY\_NOTICES.md.}
\end{itemize}
\paragraph{Preparation and review.}
Original question attributed to the cited source and UnsolvedMath contributors. Codex (GPT-6; exact serving revision not exposed) assisted with definitions, English summaries, source inspection and formatting. No new human mathematical authorship or independent proof verification is asserted. Canonical ID 30000155 and source code OWR-768-005 checked against the complete supplied registry and existing repository statements; known flagged repetitions and obvious identical conjectures excluded. Full mathematical-equivalence/literature review has not been performed. Statement preparation and final review followed contribution prompts C and D; automated checks verify structure and identity, not open status.
\label{end:30000155}

\section{Literature and context check}
Maharaj's contribution in report 54/2004, pp.~2942--2943, was read. Its function space is an internal direct sum $\bigoplus_iL(G_i)y^i$, then evaluated on one common set of rational places. Its component codes occupy the same coordinates. This is materially different from concatenating disjoint coordinate blocks. The report's tower-specific asymptotic question cannot be answered by the elementary block-direct-sum distance formula.
\section{Research attempt: Cancellation in a common-coordinate direct sum}
For disjoint coordinate blocks, $C=C_1\oplus\cdots\oplus C_r$ has minimum distance $\min_i d(C_i)$: every nonzero word has a nonzero block of weight at least that minimum, and a minimum-weight word in one block with zeros elsewhere attains it. But an internal direct sum of subspaces of the same $\mathbb F_q^N$ may have much smaller distance through cancellation.

For an explicit test over $\mathbb F_2$ with $N\geq3$, let $u=(1,\ldots,1)$ and $v=(1,\ldots,1,0)$. Their one-dimensional spans have zero intersection, because their only nonzero vectors differ, and their distances are $N$ and $N-1$. Their internal direct sum contains $u+v=(0,\ldots,0,1)$, so its minimum distance is one. The upper bound by the minimum component distance remains true because each component is a subcode; a matching lower bound fails arbitrarily badly. This is an obstruction to a proof route, not a claim that the report's particular tower has these subcodes.

A concrete evaluation model shows what additional control can supply a lower bound. Fix $N$ distinct nonzero elements $\alpha_1,\ldots,\alpha_N\in\mathbb F_q$, a positive integer $r$, and nonnegative integers $g_0,\ldots,g_{r-1}$. Consider the polynomial space
\[
 \mathcal L=\bigoplus_{i=0}^{r-1}\operatorname{span}\{Y^{rj+i}:0\leq j\leq g_i\}.
\]
All displayed exponents are distinct, so its dimension is $K=\sum_i(g_i+1)$. Put $D=\max_i(rg_i+i)$ and assume $D<N$. Evaluation at the $\alpha_j$ is injective, and the resulting code has dimension $K$ and distance at least $N-D$. Indeed a nonzero polynomial of degree at most $D$ has at most $D$ roots: division by $Y-\alpha$ at a root lowers the degree by one, proving the claim inductively. Thus no nonzero evaluation word can vanish on more than $D$ coordinates. This proof controls cancellation by a degree bound for the entire sum, not merely by component distances.

In the balanced case $g_i=g$, all exponents from zero through $D=r(g+1)-1$ occur. The distance is exactly $N-D$, since a polynomial $\prod_{j=1}^{D}(Y-\alpha_j)$ belongs to the space and vanishes at exactly those $D$ evaluation points. This verifies both the benefit and the limit of the common-degree method. It is an elementary polynomial-evaluation model; allowing $q$ to grow to accommodate unbounded $N$ does not produce a fixed-alphabet function-field tower.

The report's known bound $d\geq N-\deg G$ plays the analogous role for its full function space. To use that bound for asymptotic goodness one must prove a positive margin between $\deg G$ and $N$, together with a positive rate, along the actual selected tower construction. Good point/genus ratios of the ambient tower alone are not those missing degree estimates for these specific internal sums.
\section{Verification and adversarial audit}
The two subspaces in the cancellation example are genuinely a direct sum, so the failure is not caused by an intersection. The polynomial model proves evaluation injectivity as well as its distance bound. Its growing alphabet is explicitly separated from asymptotic goodness over a fixed field.
\paragraph{Audit corrections.} Internal sums sharing evaluation coordinates were not replaced by disjoint block sums. The asymptotic alphabet and the missing tower divisor estimates are stated explicitly.
\paragraph{External computation.} None; symbolic arguments and boundary checks only.
\section{Scope and remaining gaps}
No asymptotic estimate for the source's report-specific divisors and selected tower levels is proved. The direct-sum shortcut is refuted and a complete auxiliary evaluation calculation is supplied, but the original tower-code question remains unresolved.
\paragraph{Integration review.} Codex (GPT-6; exact serving revision not exposed) reviewed the written arguments and expanded layout and source-reference wrappers. Repository metadata and the production web renderer were checked using Python 3.13 and JavaScript with MathJax 3.2.2. This model-assisted review is not independent human proof verification; no PDF compilation or formal proof checking is asserted.

\section{Completion Estimate}
\noindent\textbf{COMPLETION: 0\%}
The calculations below the preserved source statement investigate a concrete route,
but establish no substantiated new progress toward the entire original target.
Known special cases and failed sufficient conditions are not credited as a solution.
\section{References}
The original source references and catalogue attribution are preserved in the
Sources section above. The following URLs were actually consulted for this attempt;
retention of an additional original reference does not claim fresh consultation of it.
\begin{itemize}
\item \url{https://ems.press/content/serial-article-files/45972} (consulted 2026-10-11). 
\end{itemize}
\section{Model and contributor disclosure}
OpenAI gpt-6-astra, requested reasoning configuration ultra, generated this research
attempt on 2026-10-11. The exact serving revision is not exposed. Codex GPT-6
performed orchestration, repository-format review, and written mathematical review at the declared scope; its exact serving revision
is also not exposed. The mathematical route, boundary checks and final adversarial
self-audit followed contribution prompts A then B in
\texttt{docs/contribution\_prompts.md}, after reading \texttt{CONTRIBUTING.md}.
Python 3.14.2 assembled the standalone documents; any mathematical computation is
specified in Verification. No theorem prover was used. No human mathematical
authorship, substantial human proof correction, or independent human proof checking
is asserted. Model-assisted review, source retrieval, and structural checks do not
establish correctness or novelty. Full private reasoning transcripts are unavailable
as reproducibility artifacts; the complete arguments, inputs and outcomes actually
used are stated here or in the companion batch audit.
\section{Final status}
Unresolved as a new research contribution. This is a concrete route investigation,
not a claim of new complete resolution or a current-status certification. A known
answer reproduced for context does not become a new contribution. The catalogue's
source-status label is preserved; this attempt status does not assert that a
source-labelled solved question is open.
\end{document}
